% sections/related_work.tex --- Related Work
\section{Related Work}
\label{sec:related-work}

This section places \bench{} in three nearby literatures: machine
learning for accounting and audit, graph anomaly-detection benchmarks
in the wider machine-learning community, and synthetic-data generation
for finance.  We discuss each in turn and identify the gap that
motivates the present work.

\subsection{Machine learning for accounting and audit}
\label{sec:related-work:audit-ml}

Machine-learning approaches to accounting-anomaly detection date to the
expert-system literature of the 1990s and have accelerated steadily
over the last decade.  \citet{west2016intelligent} survey the broader
landscape of intelligent financial-fraud detection and emphasise the
predominance of supervised methods on imbalanced data;
\citet{hilal2022financial} update the survey for the deep-learning
era and document a shift towards autoencoder, sequence-model, and
graph-based approaches.  Within accounting specifically,
\citet{schreyer2017detection} demonstrated that deep autoencoder
networks reconstruct routine general-ledger postings well and surface
unusual entries through reconstruction-error scoring;
\citet{schreyer2019adversarial} extend this with adversarial training
to detect synthetic ``accounting-deepfake'' postings designed to
evade autoencoder defences.  \citet{liang2021bookkeeping} apply
pattern-recognition methods to the bookkeeping graph itself, showing
that subgraph anomalies can be detected when the journal-entry table
is interpreted as a network.  Domain-specific work on
VAT-compliance violations~\citep{lahann2019vat} and on
continuous-auditing analytics~\citep{vasarhelyi2004principles}
highlights the close relationship between detection technique and the
particular accounting structure on which it operates.
\citet{guo2022picture} formalise the topology of the accounting graph
for audit and fraud-detection purposes and argue, in a position
consonant with the present paper, that accounting-graph structure
encodes information not present in the row-level table.
\citet{boersma2024audit} explores the use of complex-network methods
in audit and offers a data-driven modelling perspective that
complements the present work.

A persistent feature of this literature, however, is that empirical
evaluation is conducted on private datasets---typically a single
client's general ledger held under audit-confidentiality
constraints---and the released code, when any, does not include the
data on which the published numbers were computed.  Cross-paper
comparison is therefore not possible: a ``state-of-the-art'' method
on one corporate ledger may underperform a simpler model on another,
and the field has no standard way of detecting this from the
published numbers alone.

\subsection{Graph anomaly-detection benchmarks}
\label{sec:related-work:graph-bench}

Outside accounting, graph-based anomaly detection has converged on a
small number of public benchmarks.  The Open Graph
Benchmark~\citep{hu2020ogb} provides standardised splits and
evaluation protocols for node-, link-, and graph-level prediction
tasks across citation, biological, and product graphs and has been
widely adopted as a comparison frame for graph neural networks
(GNNs).  More recently, dedicated graph-anomaly benchmarks have
emerged: \citet{liu2024gadbench} (\textsc{GADBench}) revisit
supervised graph-anomaly detection across heterogeneous graph
families, and \citet{tang2023bond} (\textsc{BOND}) survey unsupervised
outlier-node detection on static attributed graphs.  Both benchmarks
have proved valuable in clarifying which methodological gains
generalise across graph types and which are dataset-specific.

Neither \textsc{GADBench} nor \textsc{BOND}, however, includes
accounting graphs of realistic structure.  Accounting graphs differ
from the citation, social, and product graphs that drive these
benchmarks in three ways that matter for evaluation: (i)~edges carry
\emph{signed monetary amounts} subject to a balance constraint
($\sum D = \sum C$ per posting), (ii)~the underlying business
processes (P2P, O2C, R2R) impose recurring \emph{document-flow
patterns} across edges, and (iii)~the standards-driven categorical
structure (account class, account sub-class, business process) is
itself a strong predictive signal that no general graph carries by
default.  A benchmark that exposes these features explicitly is, to
our knowledge, not available in the public-graph-benchmark space, and
it is the gap \bench{} is designed to fill.

\subsection{Synthetic data for finance}
\label{sec:related-work:synth-finance}

The third nearby literature is synthetic data generation for finance.
\citet{xu2019modeling} introduced \textsc{CTGAN} and \textsc{TVAE} for
conditional tabular generation; \citet{patki2016synthetic} earlier
presented the \textsc{Synthetic Data Vault} for hierarchical relational
data; \citet{solatorio2023realtabformer} and
\citet{kotelnikov2023tabddpm} extended tabular generation to
transformer- and diffusion-based architectures.  In the finance
domain, the \textsc{BankSim} simulator~\citep{lopez2016banksim} and
its successors have made transaction-level fraud-detection studies
reproducible and have, in particular, given rise to the widely-used
PaySim corpus.

These efforts share an important property with \bench{}: they
generate data that is independently releasable, free of confidentiality
constraints, and reproducible up to a known seed.  They differ in a
property that is decisive for accounting: \bench{} is designed not as
a tabular sampler trained from real data, but as a \emph{forward}
generator that respects the hard accounting invariants
($\sum D = \sum C$ per posting, $A = L + E$ at consolidation,
document-chain integrity for P2P/O2C flows) by
construction~\citep{ivertowski2026datasynth}.  Sample-based generators
have no native mechanism to enforce these constraints; the
forward-generation approach that \system{} implements does, and it
also exposes the underlying generative process to the user.  This
positioning is consonant with the broader argument in
\citet{cambria2021kg} that knowledge-graph-style structured generation
gives stronger reasoning purchase than purely sample-based methods.
